\documentclass[11pt]{article}

\usepackage[letterpaper,margin=1in]{geometry}
\usepackage[T1]{fontenc}
\usepackage[utf8]{inputenc}
\usepackage{lmodern}
\usepackage{microtype}
\usepackage{amsmath,amssymb,mathtools,bm}
\usepackage{enumitem}
\usepackage{xcolor}
\usepackage{booktabs}
\usepackage{array}
\usepackage{hyperref}
\usepackage[most]{tcolorbox}
\usepackage{fancyhdr}
\usepackage{titlesec}
\usepackage{parskip}
\usepackage{needspace}

\definecolor{qftblue}{HTML}{245A8D}
\definecolor{qftgreen}{HTML}{2F6B4F}
\definecolor{qftred}{HTML}{A23B3B}
\definecolor{qftgray}{HTML}{F4F6F8}

\hypersetup{
  colorlinks=true,
  linkcolor=qftblue,
  urlcolor=qftblue,
  citecolor=qftgreen,
  pdftitle={PHYS 662 Lecture 8: Lie Groups, Lie Algebras, and Representations},
  pdfauthor={Nima Lashkari}
}

\pagestyle{fancy}
\fancyhf{}
\lhead{PHYS 662: Quantum Field Theory I}
\rhead{Lecture 8}
\cfoot{\thepage}
\setlength{\headheight}{14pt}

\titleformat{\section}{\Large\bfseries\color{qftblue}}{\thesection.}{0.6em}{}
\titleformat{\subsection}{\large\bfseries\color{qftgreen}}{\thesubsection}{0.6em}{}
\titleformat{\subsubsection}{\normalsize\bfseries\color{qftred}}{\thesubsubsection}{0.6em}{}

\setlist[itemize]{leftmargin=2em,itemsep=0.25em,topsep=0.4em}
\setlist[enumerate]{leftmargin=2em,itemsep=0.25em,topsep=0.4em}

\newcommand{\dd}{\mathrm{d}}
\newcommand{\R}{\mathbb{R}}
\newcommand{\Lag}{\mathcal{L}}
\newcommand{\cM}{\mathcal{M}}
\newcommand{\SO}{\mathrm{SO}}
\newcommand{\OO}{\mathrm{O}}

\newtcolorbox{keyidea}[1][]{
  enhanced, breakable, colback=qftgray, colframe=qftblue,
  boxrule=0.8pt, arc=1.5mm, title={#1}, fonttitle=\bfseries,
  coltitle=white
}
\newtcolorbox{remarkbox}[1][]{
  enhanced, breakable, colback=white, colframe=qftgreen,
  boxrule=0.7pt, arc=1.5mm, title={#1}, fonttitle=\bfseries,
  coltitle=white
}
\newtcolorbox{warningbox}[1][]{
  enhanced, breakable, colback=white, colframe=qftred,
  boxrule=0.7pt, arc=1.5mm, title={#1}, fonttitle=\bfseries,
  coltitle=white
}


\newcommand{\Th}{\Theta}
\newcommand{\Eul}{\mathcal E}
\newcommand{\Lie}{\mathscr L}

\newcommand{\C}{\mathbb C}
\newcommand{\Z}{\mathbb Z}
\newcommand{\SU}{\mathrm{SU}}
\newcommand{\UU}{\mathrm{U}}
\newcommand{\GL}{\mathrm{GL}}
\newcommand{\SL}{\mathrm{SL}}
\newcommand{\Sp}{\mathrm{Sp}}
\newcommand{\tr}{\operatorname{tr}}
\newcommand{\Ad}{\operatorname{Ad}}
\newcommand{\ad}{\operatorname{ad}}
\newcommand{\lieg}{\mathfrak g}
\newcommand{\lieh}{\mathfrak h}
\newcommand{\su}{\mathfrak{su}}
\newcommand{\so}{\mathfrak{so}}
\newcommand{\slie}{\mathfrak{sl}}


\newcommand{\cH}{\mathcal H}
\newcommand{\Spin}{\operatorname{Spin}}
\newcommand{\id}{\mathrm{id}}
\begin{document}
\pagenumbering{gobble}
\begin{titlepage}
\centering
\vspace*{1.2in}
{\scshape\Large PHYS 662 --- Quantum Field Theory I\par}
\vspace{0.35in}
{\Huge\bfseries\color{qftblue} Lecture 8\par}
\vspace{0.15in}
{\LARGE\bfseries Lie Groups, Lie Algebras,\par}
\vspace{0.06in}
{\LARGE\bfseries and Representations\par}
\vspace{0.16in}
{\large\color{qftgreen} From Finite Transformations to Infinitesimal Structure\par}
\vspace{0.6in}
{\large Nima Lashkari\par}
{\large Fall 2026\par}
\vfill
\begin{minipage}{0.86\textwidth}
\small
A symmetry has a composition law, an action on a vector space, and
infinitesimal generators. We develop these three descriptions through
matrix groups, invariant subspaces, and Lie brackets. The adjoint action
then explains ideals and simple algebras. This gives the language needed
to construct spin and tensor representations in Lecture 9 and to study
central extensions and Noether charges in Lecture 10.

\medskip
Course page: \url{https://nimalashkari.com/quantum-field-theory-2026-part-i/}
\end{minipage}
\vspace{0.5in}
\end{titlepage}
\pagenumbering{roman}
\begingroup\small
\tableofcontents
\endgroup
\newpage
\pagenumbering{arabic}

\section*{Roadmap}
\addcontentsline{toc}{section}{Roadmap}
Lecture 7 related continuous symmetries of an action to conserved
currents. We now develop the language in which symmetry transformations
and their generators are organized. Our three questions are:
\begin{enumerate}
\item How do finite transformations compose, and what does a representation
tell us about the objects being transformed?
\item How do invariant subspaces and infinitesimal generators help us
understand a representation?
\item What distinguishes an ideal of a Lie algebra from an invariant
subspace of an arbitrary representation?
\end{enumerate}
Lecture 9 applies this language to spin, tensors, and quantum rotations.
Lecture 10 connects projective phases and central extensions to the
algebra of Noether charges.

\begin{keyidea}[Conventions]
Representation spaces are complex unless a real space is explicitly
specified. We write \(R(g)\) for a general representation and \(U(g)\)
for a unitary one. For Hermitian generators our convention is
\[
U(\theta)=e^{-i\theta^a T_a},\qquad
[T_a,T_b]=i f_{ab}{}^cT_c.
\]
The tangent matrices are \(X_a=-iT_a\), so
\([X_a,X_b]=f_{ab}{}^cX_c\). Matrix Lie algebras are described using
\(X\), with the ordinary commutator. The \(T\) convention is useful in
quantum mechanics; the \(X\) convention makes reality conditions clear.
We use \(\hbar=1\), \(\epsilon_{123}=+1\), and mostly-minus Minkowski
signature when spacetime appears.
\end{keyidea}

\section{Groups and the main matrix examples}
\label{sec:groups}

\subsection{Composition is the starting point}
A group \(G\) is a set with a multiplication satisfying closure,
associativity, an identity \(e\), and an inverse for every element:
\begin{equation}
(g_1g_2)g_3=g_1(g_2g_3),\qquad eg=ge=g,\qquad
g^{-1}g=gg^{-1}=e.
\end{equation}
Commutativity is an additional property, not an axiom. An Abelian group
obeys \(g_1g_2=g_2g_1\). Our convention is that \(g_1g_2\) acts by
applying \(g_2\) first, then \(g_1\).

A Lie group is a finite-dimensional smooth manifold for which
multiplication and inversion are smooth. Its dimension counts local real
parameters, not the size of a matrix representation. The group may have
several connected components. Its identity component \(G_0\) contains
the transformations continuously connected to \(e\).

A useful first example is a rotation in a plane:
\[
 R(\theta)=\begin{pmatrix}\cos\theta&-\sin\theta\\
 \sin\theta&\cos\theta\end{pmatrix},\qquad
 R(\theta)R(\varphi)=R(\theta+\varphi).
\]
One angle labels the transformation, even though its action on a vector
requires a two-by-two matrix. The group has dimension one and this
representation has dimension two. Keeping these two counts separate
will help with SU(2) later.

For a spacetime example, write a Poincar\'e transformation as
\(g=(\Lambda,a)\). Applying the two transformations in order gives
\[
 x\ \longmapsto\ \Lambda_2x+a_2
 \ \longmapsto\ \Lambda_1(\Lambda_2x+a_2)+a_1.
\]
Reading off the new linear part and displacement, we obtain
\begin{equation}
(\Lambda_1,a_1)(\Lambda_2,a_2)
=(\Lambda_1\Lambda_2,\,a_1+\Lambda_1a_2),\qquad
(\Lambda,a)^{-1}=(\Lambda^{-1},-\Lambda^{-1}a).
\label{eq:PoincareProduct}
\end{equation}
The Lorentz transformation rotates the translation parameter. This is
why the group is a semidirect product
\(\R^{1,d}\rtimes O(1,d)\), rather than an ordinary direct product.
The action is a functional \(S[\phi]\); its symmetry is properly written
\(S[g\cdot\phi]=S[\phi]\), with the domain and backgrounds handled
as in Lecture 7: the action assigns a number to an entire field
configuration.

\subsection{Orthogonal and unitary groups}
The real orthogonal group and its special subgroup are
\begin{equation}
O(N)=\{R\in\GL(N,\R):R^TR=I\},\qquad
SO(N)=\{R\in O(N):\det R=1\}.
\end{equation}
Orthogonality gives \((\det R)^2=1\). The determinant-minus-one
component contains improper rotations, including reflections. For
\(N\ge2\), \(SO(N)\) is the identity component. For a real multiplet, \(\phi\mapsto R\phi\) preserves
\(\phi^T\phi\), which can take any allowed value.

Similarly,
\begin{equation}
U(N)=\{U\in\GL(N,\C):U^\dagger U=I\},\qquad
SU(N)=\{U\in U(N):\det U=1\}.
\end{equation}
Unitarity implies \(|\det U|=1\). It preserves
\begin{equation}
(U\Phi)^\dagger(U\Phi)=\Phi^\dagger\Phi,
\end{equation}
not the number one unless \(\Phi\) was normalized beforehand.
For scalar field components there is usually no such normalization.

\subsection{Indefinite orthogonal and symplectic groups}
The Lorentz group satisfies
\begin{equation}
\Lambda^T\eta\Lambda=\eta,\qquad
\eta=\operatorname{diag}(1,-1,\ldots,-1).
\end{equation}
We denote it by \(O(1,d)\). Relabeling the signature convention gives
the isomorphic notation \(O(d,1)\). The condition \(\det\Lambda=1\)
defines \(SO(1,d)\), but does not impose preservation of time orientation.
The identity component is the proper orthochronous group
\(SO^+(1,d)\), with \(\Lambda^0{}_0\ge1\).

The real symplectic group preserves an antisymmetric form:
\begin{equation}
Sp(2n,\R)=\{M\in\GL(2n,\R):M^T\Omega M=\Omega\},\qquad
\Omega=\begin{pmatrix}0&I_n\\-I_n&0\end{pmatrix}.
\label{eq:symplecticGroup}
\end{equation}
It describes linear transformations preserving the canonical symplectic
form on phase space. It is different from the compact group often called
\(Sp(n)=USp(2n)\), defined by the same form-preserving condition for
complex matrices together with unitarity. Both have real dimension
\(n(2n+1)\), but their reality and compactness properties differ.
We will keep the notation explicit.

\section{Representations and invariant subspaces}

\subsection{Representations, equivalence, and faithfulness}
A group specifies how transformations compose. A representation specifies
what they act on: scalar components, vectors, tensors, or later quantum
states. A linear representation of \(G\) on \(V\) is a map into
\(\GL(V)\), the invertible linear maps on \(V\):
\begin{equation}
R:G\longrightarrow\GL(V),\qquad
\boxed{R(g_1g_2)=R(g_1)R(g_2).}
\label{eq:repDefinition}
\end{equation}
It follows that \(R(e)=I\) and \(R(g^{-1})=R(g)^{-1}\).
For Lie groups we take finite-dimensional representations to be smooth.
A representation is \emph{faithful} if its kernel is trivial: only the
identity acts as \(I\). The trivial representation sends every element
to \(1\). Thus even a non-Abelian group can have a representation with
commuting images; a representation need not retain every distinction
between group elements.

Representations \(R_1\) and \(R_2\) are equivalent if a fixed invertible
linear map \(S\) intertwines them:
\begin{equation}
R_2(g)=S R_1(g)S^{-1}\qquad\text{for all }g.
\label{eq:equivalence}
\end{equation}
For finite-dimensional unitary representations an equivalence can be
chosen unitary, in which case \(S^{-1}=S^\dagger\). A general basis
change need not be an isometry.

\subsection{Reducibility means an invariant vector subspace}
A subspace \(W\subset V\) is invariant if \(R(g)W\subset W\) for all
\(g\). A vector that starts in W stays there under every symmetry
transformation, so the components in W form a closed multiplet.
For rotations about the z-axis, the z-axis itself is an invariant line;
for the full SO(3), rotations about x or y move that line. Invariance
must always be tested against the whole group being represented.

A representation is irreducible if the only invariant subspaces
are \(0\) and \(V\). Otherwise it is reducible.
Choosing a basis beginning with a basis of a nontrivial invariant \(W\)
gives
\begin{equation}
R(g)=\begin{pmatrix}A(g)&B(g)\\0&C(g)\end{pmatrix}.
\label{eq:upperBlock}
\end{equation}
Reducibility alone gives block \emph{triangular} form. It gives a block
diagonal direct sum only if an invariant complementary subspace exists.
A representation is completely reducible if it is a direct sum of
irreducible representations.

For example, the additive group \(\R\) acts by
\begin{equation}
R(t)=\begin{pmatrix}1&t\\0&1\end{pmatrix},\qquad
R(t)R(s)=R(t+s).
\label{eq:shear}
\end{equation}
The line spanned by \(e_1\) is invariant. An independent line spanned by
\(ae_1+be_2\), \(b\ne0\), cannot be invariant for every \(t\), since
the first component changes by \(tb\). The representation is reducible
but has no invariant complement and cannot be diagonalized into two
one-dimensional representations.

\begin{warningbox}[Three different meanings of invariance]
An invariant subspace \(W\subset V\) is a property of a representation.
A normal subgroup \(H\subset G\) obeys \(gHg^{-1}=H\).
An ideal \(\lieh\subset\lieg\) obeys \([\lieg,\lieh]\subset\lieh\).
These objects live in different spaces. Irreducibility is about subspaces
of \(V\), not subgroups of \(G\).
\end{warningbox}

\subsection{Why finite-dimensional unitary representations split}
If \(W\) is invariant and \(R\) is unitary, then its orthogonal
complement is invariant too. For \(v\in W^\perp\) and \(w\in W\),
\begin{equation}
\langle w,R(g)v\rangle
=\langle R(g)^{-1}w,v\rangle=0.
\end{equation}
Hence \(V=W\oplus W^\perp\). Repeating in finite dimensions proves
complete reducibility.

Every continuous finite-dimensional representation of a compact group
can be made unitary. The idea is to average any positive inner product
over all its rotated versions, so that no group element is preferred.
For a finite group this is an ordinary average over its elements. For a
compact Lie group the sum becomes an integral with normalized invariant
Haar measure \(\dd\mu(g)\):
\begin{equation}
\langle v,w\rangle_G=\int_G\dd\mu(g)\,
\langle R(g)v,R(g)w\rangle.
\end{equation}
The result is positive and invariant. Consequently, finite-dimensional
representations of compact groups decompose into irreducible pieces.
There is also a complete-reducibility theorem for finite-dimensional
representations of semisimple Lie algebras over \(\R\) or \(\C\);
it does not require the representation to be unitary.

In infinite-dimensional Hilbert spaces, continuous spectra can require
direct integrals and additional analytic hypotheses. Our spin and tensor
decompositions are finite dimensional, where the argument above applies.

\subsection{Schur's lemma: what commutes with an irreducible action?}
For a finite-dimensional complex irreducible representation, any linear
operator \(A\) commuting with all \(R(g)\) is a scalar multiple of the
identity. To see this, choose an eigenvalue \(a\); the nonzero eigenspace
\(\ker(A-aI)\) is invariant because
\(A R(g)v=R(g)Av=aR(g)v\) whenever \(Av=av\). Irreducibility makes
that eigenspace all of \(V\), so \(A=aI\).
This is Schur's lemma. In Lecture 9 it will let us label spin
representations by the scalar value of an operator commuting with all
generators.

\Needspace{9\baselineskip}
\Needspace{10\baselineskip}
\section{From a Lie group to its Lie algebra}
\label{sec:algebra}
\subsection{One-parameter subgroups and generators}
A differentiable one-parameter matrix subgroup satisfies
\(U(s+t)=U(s)U(t)\) and \(U(0)=I\). Differentiating with respect to \(t\)
at zero gives
\[
\frac{\dd U(s)}{\dd s}=U(s)X,
\qquad X=\left.\frac{\dd U(t)}{\dd t}\right|_{t=0},
\qquad U(s)=e^{sX}.
\]
For the planar rotation above, the derivative is especially concrete:
\[
 X=R'(0)=\begin{pmatrix}0&-1\\1&0\end{pmatrix},\qquad X^2=-I.
\]
The even and odd terms of \(e^{\theta X}\) sum to
\(I\cos\theta+X\sin\theta=R(\theta)\). Thus the generator is the
first-order change; exponentiation builds the finite motion from it.
The tangent vectors \(X\) at the identity form the Lie algebra
\(\lieg\). For a unitary representation, differentiating
\(U(t)^\dagger U(t)=I\) shows that \(X^\dagger=-X\).
Physicists usually write \(X=-iT\), with \(T\) Hermitian:
\begin{equation}
 U(\theta)=e^{-i\theta^aT_a},
 \qquad [T_a,T_b]=i f_{ab}{}^cT_c.
 \label{eq:generator-convention}
\end{equation}
Here the parameters and structure constants are real for a real Lie
algebra in a real basis. In the anti-Hermitian basis \(X_a=-iT_a\),
\begin{equation}
 [X_a,X_b]=f_{ab}{}^cX_c.
 \label{eq:antihermitian}
\end{equation}
The Lie algebra itself is independent of either convention. For a
nonunitary representation the matrices \(T_a\) need not be Hermitian;
for a nonfaithful representation some generators may act as zero.

The bracket is bilinear, antisymmetric, and obeys the Jacobi identity,
\begin{equation}
 [X,[Y,Z]]+[Y,[Z,X]]+[Z,[X,Y]]=0.
 \label{eq:jacobi}
\end{equation}
For matrices this follows by expanding the twelve products and cancelling
them in pairs. In components it constrains the structure constants:
\[
 f_{bc}{}^d f_{ad}{}^e+
 f_{ca}{}^d f_{bd}{}^e+
 f_{ab}{}^d f_{cd}{}^e=0.
\]
A Lie algebra representation is a linear map \(\rho\) satisfying
\(\rho([X,Y])=[\rho(X),\rho(Y)]\). Differentiating a group
representation always produces such a map.

\subsection{Why the commutator measures noncommutativity}
Compare the two orders of small transformations generated by \(X\) and
\(Y\) by forming their group commutator. Multiplication of its four power
series gives
\begin{equation}
 e^{sX}e^{tY}e^{-sX}e^{-tY}
 =I+st[X,Y]+O(s^2t,st^2).
 \label{eq:group-commutator}
\end{equation}
The result is not generally the identity. The surviving term is another
tangent direction, which explains closure of the Lie algebra under
commutators. In the Hermitian convention this term is
\(-ist f_{ab}{}^cT_c\) when \(X=-iT_a\), \(Y=-iT_b\).

Two useful forms of the same local information are
\begin{align}
 e^Ae^B&=\exp\!\left(A+B+\tfrac12[A,B]+\cdots\right),
 \label{eq:bch}\\
 e^AYe^{-A}&=Y+[A,Y]+\tfrac12[A,[A,Y]]+\cdots.
 \label{eq:conjugation}
\end{align}
The first is the Baker--Campbell--Hausdorff formula, understood here near
the identity. The second follows by differentiating
\(e^{tA}Ye^{-tA}\) and is the expansion of conjugation.

\subsection{Extracting the bracket from a composition law}
We can now see where the structure constants enter the group law itself.
The calculation below keeps the mixed second-order terms: first-order
terms alone cannot distinguish the two orders of composition. Let \(g(\theta)g(\varphi)=g(F(\theta,\varphi))\), with
\(g(0)=e\). The identity conditions imply
\begin{equation}
 F^c(\theta,\varphi)=\theta^c+\varphi^c+
 B_{ab}{}^c\theta^a\varphi^b+O(3).
 \label{eq:coordinate-composition}
\end{equation}
There are no pure quadratic terms in \(\theta\) or \(\varphi\), but the
mixed coefficients need not be antisymmetric in arbitrary coordinates.
Expand a representation as
\[
 R(g(\theta))=I-i\theta^aT_a+
 \tfrac12\theta^a\theta^b C_{ab}+O(\theta^3),
 \qquad C_{ab}=C_{ba}.
\]
Comparing mixed terms in
\(R(g(\theta))R(g(\varphi))=R(g(F(\theta,\varphi)))\) yields
\[
 -T_aT_b=C_{ab}-iB_{ab}{}^cT_c.
\]
Subtract the equation with \(a\) and \(b\) exchanged. The symmetric
second derivative drops out, leaving
\begin{equation}
 [T_a,T_b]=i\bigl(B_{ab}{}^c-B_{ba}{}^c\bigr)T_c,
 \qquad f_{ab}{}^c=B_{ab}{}^c-B_{ba}{}^c.
 \label{eq:structure-from-composition}
\end{equation}
Thus the bracket detects precisely the antisymmetric part of local
composition. This computation may be performed in a faithful local
realization; other representations inherit the same constants, although
their images may not distinguish every generator.

In exponential coordinates, \(g(\theta)=\exp(\theta^aX_a)\),
Eq.~\eqref{eq:bch} gives
\[
 B_{ab}{}^c=\tfrac12 f_{ab}{}^c,
 \qquad C_{ab}=-\tfrac12\{T_a,T_b\}.
\]
In particular,
\begin{equation}
 e^{-i\theta^aT_a}
 =I-i\theta^aT_a-\tfrac12\theta^a\theta^bT_aT_b+O(\theta^3).
 \label{eq:second-order}
\end{equation}
The quadratic term has a minus sign. Also,
\(g(\theta)^{-1}=g(-\theta)\) holds in these coordinates, but need not
hold in an arbitrary parametrization.

\subsection{Algebras of the matrix groups}
The infinitesimal constraints follow by putting \(M=I+tX+O(t^2)\)
into each defining group equation. For example,
\[
 M^TM=I+t(X^T+X)+O(t^2)=I
 \quad\Longrightarrow\quad X^T=-X.
\]
A real antisymmetric matrix has one independent entry for each pair
\(i<j\), giving \(N(N-1)/2\) rotation generators. The same procedure
produces the other rows below. All dimensions are real.
\begin{center}
\small
\begin{tabular}{@{}lll@{}}
\toprule
Group & Constraint on \(X\) & Dimension \\
\midrule
\(\UU(N)\) & \(X^\dagger+X=0\) & \(N^2\) \\
\(\SU(N)\) & \(X^\dagger+X=0,\ \tr X=0\) & \(N^2-1\) \\
\(\OO(N),\ \SO(N)\) & \(X^T+X=0\) & \(N(N-1)/2\) \\
\(\OO(1,d),\ \SO^+(1,d)\) & \(X^T\eta+\eta X=0\) & \(d(d+1)/2\) \\
\(\Sp(2n,\R)\) & \(X^T\Omega+\Omega X=0\) & \(n(2n+1)\) \\
\bottomrule
\end{tabular}
\end{center}
For example, differentiating \(\det e^{tX}=e^{t\tr X}\) explains the
traceless condition for a special unitary matrix. A real antisymmetric
matrix is already traceless, so imposing determinant one does not
remove an additional generator from \(\OO(N)\). The groups
\(\OO(N)\) and \(\SO(N)\) have the same Lie algebra because their
identity components agree.

For the symplectic group, write \(X\) in \(n\)-by-\(n\) blocks.
The constraint is equivalent to
\begin{equation}
 X=\begin{pmatrix}A&B\\ C&-A^T\end{pmatrix},
 \qquad B=B^T,\quad C=C^T.
 \label{eq:symplectic-algebra}
\end{equation}
There are \(n^2+n(n+1)=n(2n+1)\) independent entries.
It is \(\Omega X\), not generally \(X\), that is symmetric.
The compact group usually denoted \(\Sp(n)\) has the same real dimension,
but is a different real group: in its complex \(2n\)-dimensional defining
representation one imposes both unitarity and preservation of the
symplectic form.

\subsection{From local generators to a connected group}
The exponential map covers a neighborhood of the identity. Every element
of a connected Lie group is a finite product of exponentials, though a
single exponential need not suffice for a general connected group.
An invariant subspace for all generators is therefore invariant under
the connected group, and conversely.

The algebra determines local composition, but not global identifications
or disconnected components. If its bracket vanishes, the identity
component is abelian; the full group need not be. For example, a
reflection in \(\OO(2)\) conjugates a rotation through \(\theta\) into
one through \(-\theta\). Lecture 9 develops the global question using
covering groups and projective representations.

\Needspace{10\baselineskip}
\section{Ideals, simple algebras, and real forms}
\label{sec:ideals}
The algebra has an action on itself: conjugation moves one generator
into a linear combination of generators. Studying its invariant
subspaces reveals the internal structure of the symmetry algebra.

\subsection{Subalgebras, ideals, and the adjoint representation}
A \emph{subalgebra} \(\lieh\subseteq\lieg\) is a linear subspace closed
under its own bracket: \([\lieh,\lieh]\subseteq\lieh\).
An \emph{ideal} satisfies the stronger condition
\begin{equation}
 [\lieg,\lieh]\subseteq\lieh.
 \label{eq:ideal}
\end{equation}
This means that bracketing with any generator of the full algebra keeps
one inside \(\lieh\). It is the infinitesimal version of a normal subgroup,
whose defining condition is \(gHg^{-1}=H\) for every \(g\in G\).

The connection is made by the adjoint representation. For a matrix group,
\begin{equation}
 \Ad_g(X)=gXg^{-1},\qquad \ad_X(Y)=[X,Y].
 \label{eq:adjoint}
\end{equation}
Conjugation preserves the Lie algebra, and
\(\Ad_{g_1g_2}=\Ad_{g_1}\Ad_{g_2}\). Differentiating gives \(\ad\);
the Jacobi identity says
\([\ad_X,\ad_Y]=\ad_{[X,Y]}\), so it too is a representation.
Ideals are exactly invariant subspaces of the adjoint Lie algebra
representation. For connected \(G\), this also means invariance under
\(\Ad_g\) for all \(g\). The connected subgroup generated by the
exponentials of an ideal is normal in a connected group; it need not be
closed, so global subgroup statements require care.

For example, take rotation generators obeying
\([J_i,J_j]=i\epsilon_{ijk}J_k\). The span of \(J_3\) is a subalgebra,
but not an ideal, since \([J_1,J_3]=-iJ_2\).
By contrast, translations form an abelian ideal of the Poincar\'e algebra:
bracketing a Lorentz generator with a translation gives a translation.
This is the infinitesimal origin of the semidirect product in
Eq.~\eqref{eq:PoincareProduct}.

\subsection{Simple and semisimple algebras}
A Lie algebra is \emph{simple} if it is nonabelian and has no nonzero
proper ideals. It is \emph{semisimple} if it has no nonzero solvable
ideals.\footnote{An algebra is solvable if its derived series
\(\lieh^{(0)}=\lieh\), \(\lieh^{(r+1)}=[\lieh^{(r)},\lieh^{(r)}]\)
eventually becomes zero.} In finite dimensions over \(\R\) or \(\C\),
this is equivalent
to being a direct sum of simple Lie algebras; it is also equivalent to
having no nonzero abelian ideals.

Neither condition says that every pair of independent generators has
a nonzero commutator. The two independent traceless diagonal Hermitian
matrices in \(\su(3)\) commute. Generators belonging to different simple
factors of a semisimple direct sum also commute.

What is true is
\begin{equation}
 [\lieg,\lieg]=\lieg
 \quad\text{for a semisimple Lie algebra},
 \label{eq:perfect}
\end{equation}
where \([\lieg,\lieg]\) denotes the linear span of all commutators.
For a simple algebra that span is a nonzero ideal and hence is the whole
algebra; the result then follows factor by factor. Thus, in any
finite-dimensional representation,
\begin{equation}
 \tr\rho(X)=0\qquad\text{for every }X\in\lieg,
 \label{eq:semisimple-trace}
\end{equation}
since every \(X\) is a sum of commutators and the trace of a commutator
vanishes. Thus tracelessness follows from the structure of the algebra.

Some useful examples, stated at the level of real Lie algebras, are
\begin{itemize}
\item \(\su(N)\) is simple for \(N\geq2\).
\item \(\so(3)\) and \(\so(N)\) for \(N\geq5\) are simple;
\(\so(4)\cong\su(2)\oplus\su(2)\) is semisimple but not simple,
and \(\so(2)\) is abelian.
\item \(\mathfrak u(N)=\su(N)\oplus\mathfrak u(1)\) is reductive,
meaning a direct sum of a semisimple algebra and an abelian center;
it is not semisimple because of its central \(\mathfrak u(1)\) ideal.
\item \(\so(1,d)\) is real simple for \(d\geq2\), while
\(\so(1,1)\) is one-dimensional and abelian.
\item The Poincar\'e algebra is not semisimple because its translation
ideal is nonzero and abelian.
\end{itemize}

A simple Lie algebra need not belong to a group with no nontrivial
normal subgroups. For example, \(\SU(N)\) has the finite center
\(\Z_N\), although its Lie algebra is simple. Discrete central
subgroups have no tangent directions and therefore are not detected by
the Lie algebra. Likewise the algebra direct sum for \(\mathfrak u(N)\)
does not give a global direct product without a quotient:
\[
 \UU(N)\cong\bigl(\SU(N)\times\UU(1)\bigr)/\Z_N.
\]
The map \((S,z)\mapsto zS\) is onto; its kernel consists of
\((z^{-1}I,z)\) with \(z^N=1\).

\Needspace{34\baselineskip}
\subsection{Real algebras and their complexifications}
This subsection records distinctions needed later for Lorentz
representations. Its main point is that allowing complex coefficients
can identify algebras that were different over the reals. The detailed
Lorentz applications can be revisited when we study relativistic fields.
The scalars over which an algebra is defined matter. For a real algebra,
its complexification \(\lieg_{\C}=\lieg\otimes_{\R}\C\) allows complex
linear combinations of its generators. Distinct real algebras can have
the same complexification:
\begin{equation}
 \su(2)_{\C}\cong\slie(2,\C),\qquad
 \slie(2,\R)_{\C}\cong\slie(2,\C).
 \label{eq:real-forms}
\end{equation}
Nevertheless \(\su(2)\) and \(\slie(2,\R)\) are not isomorphic as real
Lie algebras. One is the compact real form; the other is noncompact.
The correct real isomorphisms relevant here are
\[
 \su(2)\cong\so(3),\qquad
 \slie(2,\R)\cong\so(1,2).
\]
At the connected group level,
\(\SO^+(1,2)\cong\mathrm{PSL}(2,\R)
=\SL(2,\R)/\{I,-I\}\).

Similarly, the Euclidean rotation algebra and Lorentz algebra in four
dimensions must not be interchanged:
\begin{align}
 \so(4)&\cong\su(2)\oplus\su(2)
 &&\text{as real algebras},\\
 \so(1,3)_{\C}&\cong\slie(2,\C)\oplus\slie(2,\C)
 &&\text{as complex algebras}.
 \label{eq:lorentz-complexification}
\end{align}
The real algebra \(\so(1,3)\) is simple even though its
complexification splits. The familiar two spin labels of complex
finite-dimensional Lorentz representations come from the second line.
Their derivation belongs to the later discussion of relativistic fields.

\begin{warningbox}[Algebra classification is not the whole group problem]
Structure constants in a specified basis determine a Lie algebra.
They do not determine a representation or a unique global Lie group.
Classifying semisimple algebras is a central part of the subject, but
more general Lie algebras also contain solvable ideals and extensions.
Global quotients and disconnected components are additional group data.
\end{warningbox}

\clearpage
\section*{Questions and checks}
\addcontentsline{toc}{section}{Questions and checks}
\begin{enumerate}
\item Derive the Poincar\'e composition law by applying two transformations
to \(x^\mu\). Why is the product semidirect?
\item For the shear representation \(R(t)=\left(\begin{smallmatrix}1&t\\0&1\end{smallmatrix}\right)\),
identify an invariant line and prove that it has no invariant complement.
Which hypothesis of the unitary complete-reducibility argument fails?
\item Expand the group commutator in Eq.~\eqref{eq:group-commutator}
through order \(st\). Track the sign after putting \(X=-iT_a\).
\item Derive the infinitesimal conditions for \(\SO(N)\), \(\SU(N)\),
and \(\Sp(2n,\R)\), and count their independent real parameters.
\item Give a subalgebra of \(\su(2)\) that is not an ideal, and two
independent commuting generators of \(\su(3)\). Why is neither example
in conflict with simplicity?
\item Distinguish these four statements: a group is connected; an algebra
is simple; a representation is irreducible; a representation is faithful.
Use the shear representation, the trivial representation of \(\SU(2)\),
and the algebra \(\su(2)\) for examples.
\item Use the Jacobi identity to prove that the span \([\lieg,\lieg]\)
of commutators is an ideal. Explain why it equals \(\lieg\) for a
semisimple algebra, and deduce tracelessness in a finite-dimensional
representation.
\item Distinguish a real algebra from its complexification using
\(\su(2)\) and \(\slie(2,\R)\). Why does the complex splitting of
\(\so(1,3)_{\C}\) not contradict real simplicity of \(\so(1,3)\)?
\end{enumerate}

\clearpage
\section*{Summary}
\addcontentsline{toc}{section}{Summary}
\begin{keyidea}[Three levels of a symmetry]
The group describes finite composition. A representation specifies the
space on which it acts. The Lie algebra records the local generators
and their commutators. Keeping these levels distinct is the starting
point for the spin and tensor constructions in Lecture 9.
\end{keyidea}
\begin{itemize}
\item A Lie group's dimension counts continuous parameters, not the
dimension of a chosen representation.
\item Irreducibility concerns invariant vector subspaces. Finite-dimensional
unitary representations are completely reducible; compact groups allow
an invariant positive inner product to be constructed by averaging.
\item Schur's lemma makes an operator commuting with an irreducible
complex action a scalar. This will determine the spin Casimir.
\item Differentiation gives tangent matrices \(X_a=-iT_a\), with
\([X_a,X_b]=f_{ab}{}^cX_c\). Antisymmetric local composition determines
the structure constants.
\item Ideals are invariant subspaces of the adjoint action. Simple and
semisimple describe algebras, independently of irreducibility of some
other representation. A semisimple algebra is spanned by commutators.
\item Reality conditions, global quotients, and disconnected components
are additional data. A Lie algebra does not determine a unique global
group; Lecture 9 makes this explicit for rotations.
\end{itemize}

\section*{Further reading}
\addcontentsline{toc}{section}{Further reading}
These notes reorganize the supplied symmetry and representation slide
sets into Lectures 8--10. For mathematical background and further proofs:
\begin{itemize}
\item B. C. Hall, \emph{An Elementary Introduction to Groups and
Representations}, \href{https://arxiv.org/abs/math-ph/0005032}{arXiv:math-ph/0005032}.
\item P. Etingof, \emph{Lie Groups and Lie Algebras I}, MIT 18.745,
\href{https://ocw.mit.edu/courses/18-745-lie-groups-and-lie-algebras-i-fall-2020/mit18_745_f20_lec18.pdf}{Lecture 18: representations of semisimple Lie algebras}.
\end{itemize}
\end{document}
